% Generated draft: requires mathematical and source review.
\documentclass[11pt]{article}
\usepackage[margin=25mm]{geometry}
\usepackage{fontspec}
\setmainfont{Latin Modern Roman}
\usepackage{amsmath,amssymb,mathtools,unicode-math}
\setmathfont{Latin Modern Math}
\usepackage{enumitem,needspace,xurl,hyperref}
\setlength{\parindent}{0pt}
\setlength{\parskip}{5pt}
\setlength{\emergencystretch}{3em}
\newcommand{\N}{\mathbb{N}}
\newcommand{\Z}{\mathbb{Z}}
\newcommand{\Q}{\mathbb{Q}}
\newcommand{\R}{\mathbb{R}}
\newcommand{\C}{\mathbb{C}}
\newcommand{\F}{\mathbb{F}}
\newcommand{\A}{\mathbb{A}}
\newcommand{\PP}{\mathbb{P}}
\newcommand{\E}{\mathbb{E}}
\newcommand{\eps}{\varepsilon}
\newcommand{\dd}{\,\mathrm{d}}
\begin{document}
\section*{UnsolvedMath ID: 30004596}
\textbf{Source problem code:} catalogue entry
\subsection{Definitions and mathematical statement}
Let \((X,\mathcal X)\) denote the underlying measurable space and let \(\lambda\) be an intensity measure for which the Poisson process in the source is defined. Write \(\eta\) for that Poisson random counting measure: for measurable \(B\) with \(\lambda(B)<\infty\), \(\eta(B)\) has the Poisson distribution with mean \(\lambda(B)\), and counts on disjoint measurable sets are independent. Let \(\eta'\) be an independent copy of \(\eta\). For a measurable real-valued function \(f\) of a counting configuration \(\mu\), define the add-one difference by
\[
D_xf(\mu):=f(\mu+\delta_x)-f(\mu),
\]
where \(\delta_x\) is the unit point mass at \(x\). The configuration space, its measurable structure, and the assumptions ensuring that these quantities and the following integrals are defined are not specified in the supplied record.

The source asks whether the Poisson OSSS-type inequality extends to an arbitrary stopping set \(Z\), beyond stopping sets expressible as unions of increasing families of stopping sets satisfying continuity assumptions. Its proposed inequality is
\[
\mathbb E\bigl[|f(\eta)-f(\eta')|\bigr]
\leq
2\int_X \mathbb P(x\in Z)\,
\mathbb E\bigl[|D_xf(\eta)|\bigr]\,\lambda(\mathrm dx).
\tag{1}
\]
For measurable \(f\) taking values in \(\{0,1\}\), the proposed formulation is
\[
\operatorname{Var}(f(\eta))
\leq
\int_X \mathbb P(x\in Z)\,
\mathbb E\bigl[|D_xf(\eta)|\bigr]\,\lambda(\mathrm dx).
\tag{2}
\]
For this Boolean-valued class, independence and equality in distribution give
\[
\mathbb E\bigl[|f(\eta)-f(\eta')|\bigr]
=2\operatorname{Var}(f(\eta)),
\]
so (1) and (2) are equivalent within that class. This does not establish equivalence between the Boolean assertion and (1) for general real-valued functions.

\textbf{Unresolved hypotheses.} The record does not define a stopping set, specify the admissible underlying spaces or test sets used in that definition, or say whether auxiliary randomization is allowed. In particular, it does not state a relationship requiring \(Z\) to determine \(f(\eta)\). Without a restriction linking \(Z\) and \(f\), the displayed assertion cannot hold for arbitrary pairs: allowing \(Z=\varnothing\) makes the right-hand side zero even when \(f(\eta)\) is nonconstant. A possible intended requirement is measurability of \(f(\eta)\) with respect to the information revealed on \(Z\), together with any permitted auxiliary randomization, but the precise requirement is absent and is not adopted here as a supplied hypothesis. The continuity assumptions and the definition of the increasing family in the restricted formulation are also absent. Accordingly, (1)--(2) record the proposed target, whose admissible class still requires clarification; they are not a complete universally quantified conjecture as supplied.
\subsection{Short English statement}
The proposed problem asks whether a Poisson OSSS-type inequality extends from stopping sets obtained through an increasing family with continuity assumptions to arbitrary stopping sets. The supplied formulation omits the definition of a stopping set and the necessary relationship between the revealed information and \(f\), so its precise admissible class remains unresolved. The variance form is equivalent to the absolute-difference form for \(\{0,1\}\)-valued functions. The supplied open status is dated curation, not an independently verified conclusion.
\subsection{Sources}
\begin{itemize}
\item Supplied local registry record: permanent UnsolvedMath ID \(30004596\), title \emph{Poisson Variance Inequalities for Arbitrary Stopping Sets}, collection \emph{Oberwolfach Reports Open Problems}, exact source problem code \texttt{catalogue entry}. The record gives proposed year 2021 and does not name a proposer.
\item Supplied source citation: \emph{Spatial Networks and Percolation (2022)}, Oberwolfach Reports, pp.~149--224, DOI \texttt{10.4171/owr/2021/4}. \url{https://doi.org/10.4171/owr/2021/4}. The supplied background describes Oberwolfach Reports as reports documenting mathematical workshops and their open-problem discussions.
\item Supplied literature source, labeled ``Primary paper on Poisson OSSS inequalities'': \url{https://arxiv.org/abs/2101.07180}. No authors, paper title, or theorem locator are supplied.
\item Supplied literature source, labeled ``Oberwolfach Report 4/2021'': \url{https://doi.org/10.4171/owr/2021/4}.
\item Literature assessment dated exactly \texttt{2026-08-22}: the supplied curation labels the problem open and states that the Poisson randomized-stopping-set variance inequality is proved only in restricted settings, while validity for an arbitrary stopping set is still stated as open in the primary account. This is supplied dated curation derived from the OWR record, not independently verified here; its time-sensitive status and cited claims require independent verification.
\item The registry assigns the statement-status label \texttt{corrected\_verified} and reports that the repaired standalone statement was checked against the official Oberwolfach report. This is a supplied verification claim, not verification performed for this draft, and it does not resolve the omitted hypotheses identified above.
\end{itemize}
\end{document}
